% ---------------------------------------------------------------- tipografia
\usepackage{microtype}
\usepackage{xurl}          % parte las URL largas por cualquier sitio razonable
\usepackage{needspace}     % evita partir un diagrama en dos paginas
\usepackage{graphicx}

% Bloques de codigo: sin esto, una linea larga se sale del margen en vez de
% partirse, y varias ordenes del manual pasan de 80 columnas.
\usepackage{fvextra}
\DefineVerbatimEnvironment{Highlighting}{Verbatim}{%
  breaklines,breakanywhere,commandchars=\\\{\},fontsize=\small}
\fvset{breaklines=true,breakanywhere=true}

% Fondo suave para el codigo, para separarlo del cuerpo de un vistazo.
\usepackage{tcolorbox}
\tcbuselibrary{skins,breakable}
\definecolor{fondocodigo}{HTML}{F6F6F3}
\definecolor{bordecodigo}{HTML}{D9D9D2}
\renewenvironment{Shaded}{%
  \begin{tcolorbox}[breakable,enhanced,boxrule=0.4pt,arc=2pt,
    left=6pt,right=6pt,top=5pt,bottom=5pt,
    colback=fondocodigo,colframe=bordecodigo]}{\end{tcolorbox}}

% ------------------------------------------------------------------- tablas
% El manual tiene tablas de tres columnas con celdas largas; a cuerpo normal se
% desbordan.
\usepackage{etoolbox}
\AtBeginEnvironment{longtable}{\small}
\setlength{\LTpre}{8pt}
\setlength{\LTpost}{10pt}
\renewcommand{\arraystretch}{1.18}

% --------------------------------------------------------------- encabezados
\usepackage{titlesec}
\definecolor{acento}{HTML}{1F4E5F}
\titleformat{\chapter}[display]
  {\normalfont\sffamily\huge\bfseries\color{acento}}
  {\sffamily\normalsize\mdseries\MakeUppercase{\chaptertitlename\ \thechapter}}
  {6pt}{\Huge}
\titlespacing*{\chapter}{0pt}{0pt}{28pt}
\titleformat{\section}{\normalfont\sffamily\Large\bfseries\color{acento}}{\thesection}{0.7em}{}
\titleformat{\subsection}{\normalfont\sffamily\large\bfseries}{\thesubsection}{0.7em}{}
\titleformat{\subsubsection}{\normalfont\sffamily\normalsize\bfseries}{\thesubsubsection}{0.7em}{}

% ------------------------------------------------------------ pie y cabecera
\usepackage{fancyhdr}
\pagestyle{fancy}
\fancyhf{}
\fancyhead[LE,RO]{\sffamily\small\thepage}
\fancyhead[RE]{\sffamily\small\nouppercase{\leftmark}}
\fancyhead[LO]{\sffamily\small\nouppercase{\rightmark}}
\renewcommand{\headrulewidth}{0.3pt}
\fancypagestyle{plain}{\fancyhf{}\renewcommand{\headrulewidth}{0pt}%
  \fancyfoot[C]{\sffamily\small\thepage}}

% ------------------------------------------------------------------- citas
% Los avisos del manual van como cita; se marcan con una barra lateral.
\usepackage{mdframed}
\definecolor{bordeaviso}{HTML}{1F4E5F}
\renewenvironment{quote}{%
  \begin{mdframed}[leftline=true,rightline=false,topline=false,bottomline=false,
    linewidth=2pt,linecolor=bordeaviso,leftmargin=0pt,innerleftmargin=10pt,
    innertopmargin=4pt,innerbottommargin=4pt,skipabove=8pt,skipbelow=8pt]}%
  {\end{mdframed}}

% Los colores de los enlaces se fijan cuando hyperref ya esta cargado.
\AtBeginDocument{\hypersetup{linkcolor=acento,urlcolor=acento,citecolor=acento,
  linktoc=all,pdftitle={HACU — documentación técnica},
  pdfauthor={AudacIA · Universidad Simón Bolívar}}}

\usepackage[none]{hyphenat}
\sloppy


\newcommand{\ManualVersion}{Versión de 17 de septiembre de 2026 · suite de regresión 517/517}

% -------------------------------------------------------------------- portada
\usepackage{tikz}
\renewcommand{\maketitle}{%
  \begin{titlepage}
  \thispagestyle{empty}
  \begin{tikzpicture}[remember picture,overlay]
    \fill[acento] (current page.north west) rectangle
      ([yshift=-6.2cm]current page.north east);
    \node[anchor=north west,text=white,inner sep=0pt]
      at ([shift={(2.6cm,-2.0cm)}]current page.north west)
      {\sffamily\bfseries\fontsize{58}{62}\selectfont HACU};
    \node[anchor=north west,text=white,inner sep=0pt,align=left]
      at ([shift={(2.65cm,-3.7cm)}]current page.north west)
      {\sffamily\fontsize{13}{18}\selectfont Hardware de Audacia\\
       \sffamily\fontsize{13}{18}\selectfont de Comunicación Universitaria};
  \end{tikzpicture}%
  \par
  \vspace*{7.4cm}
  \begin{flushleft}
    {\sffamily\LARGE\bfseries\color{acento}
      Documentación técnica\\[5pt]y guía de replicación\par}
    \vspace{14pt}
    {\color{acento}\rule{5cm}{2pt}\par}
    \vspace{18pt}
    \parbox{12.2cm}{\raggedright
      Asistente conversacional local para la exhibición de AudacIA: modelo de
      lenguaje, recuperación sobre el corpus institucional, voz y ventana de
      exhibición. Todo en la máquina, sin llamadas externas.\par}
  \end{flushleft}
  \vfill
  \begin{flushleft}
    {\sffamily\large AudacIA · Centro de Excelencia en IA y Robótica\par}
    \vspace{3pt}
    {\sffamily Universidad Simón Bolívar · Barranquilla\par}
    \vspace{12pt}
    {\sffamily\small\color{acento!75!black}\ManualVersion\par}
  \end{flushleft}
  \end{titlepage}
}
